\documentclass[10pt,a4paper]{article}
\usepackage[margin=2.1cm]{geometry}
\usepackage{amsmath,amssymb,booktabs}
\usepackage[T1]{fontenc}
\usepackage{lmodern}
\usepackage[hidelinks]{hyperref}
\setlength{\parindent}{0pt}
\setlength{\parskip}{5pt}
\newcommand{\code}[1]{\texttt{\detokenize{#1}}}
\newcommand{\mueff}{\mu_\mu^{\mathrm{r}}}
\begin{document}
\begin{center}
{\Large\bfseries Pair processes and muon decay in GRACE}\\[4pt]
{\small Implementation summary, 10 November 2026, commit \code{1b78c044}}
\end{center}
This note describes \code{pair_treatment: evolved}, optimized \code{spectral}
decay, and \code{grey_decay} (the code spelling of ``decay\_grey'').
Equations are local fluid-frame collision equations unless stated otherwise.
Energies, temperature and chemical potentials are in MeV, with $k_B=1$.
Charged-lepton chemical potentials include rest mass. A dot denotes local proper time.

\section{Evolved thermal pair processes}
The explicit heavy-flavour channels are
\begin{equation}
 e^-+e^+\rightleftharpoons\nu_\ell+\bar\nu_\ell,\qquad
 NN\rightleftharpoons NN+\nu_\ell+\bar\nu_\ell,\qquad
 \gamma^*\rightleftharpoons\nu_\ell+\bar\nu_\ell .
\end{equation}
Here $\ell=\mu,\tau$. The plasmon model is optional. These reactions produce
neutrino pairs and do not require charged muons. The electron-pair and
bremsstrahlung kernels come from BNS\_NURATES. The former uses a relativistic,
massless-electron approximation. The optional plasmon model retains transverse modes.

For one neutrino helicity, let $w_i$ include the quadrature measure
$d\Pi=\epsilon^2d\epsilon/[2\pi^2(\hbar c)^3]$.
At every implicit trial, the two evolved moments determine an FD spectral fit:
\begin{equation}
 f_{s,i}=\frac{1}{\exp(\epsilon_i/T_{\nu,s}-\eta_s)+1},\qquad
 n_s=\sum_i w_i f_{s,i},\qquad J_s=\sum_i w_i\epsilon_i f_{s,i}.
 \label{eq:fit}
\end{equation}
$T_{\nu,s}$ and $\eta_s$ are fitted independently for each species. They need
not equal the matter values. This is a spectral closure of grey M1 moments.
For production and absorption kernels $R^p,R^a$ in $\mathrm{cm^3\,s^{-1}}$,
\begin{align}
 R^p_{ij}&=e^{-(\epsilon_i+\epsilon_j)/T}R^a_{ij},\\
 I_{ij}&=w_iw_j\left[R^p_{ij}(1-f_i)(1-\bar f_j)-R^a_{ij}f_i\bar f_j\right],\\
 \dot n_\nu=\dot n_{\bar\nu}&=r=\sum_{ij}I_{ij},\qquad
 Q_\nu=\sum_{ij}\epsilon_i I_{ij},\qquad
 Q_{\bar\nu}=\sum_{ij}\epsilon_j I_{ij}.
 \label{eq:event}
\end{align}
Both partners therefore receive the same number increment. Pair reactions
preserve their lepton-number difference and give no direct change in $Y_\mu$.
For the tau aggregate, GRACE divides the moments by multiplicity two before
the fit and restores that factor in the sources.

The damping includes own-species Pauli blocking:
\begin{equation}
 j_i=\sum_jw_jR^p_{ij}(1-\bar f_j),\qquad
 a_i=\sum_jw_jR^a_{ij}\bar f_j,\qquad
 C_i=j_i-(j_i+a_i)f_i.
\end{equation}
With $L_E=\sum_iw_i\epsilon_i(j_i+a_i)f_i$, the isotropic-kernel M1 source is
\begin{equation}
 G^\alpha=Q u^\alpha-\lambda_E H^\alpha,\qquad \lambda_E=L_E/J.
\end{equation}
$H^\alpha$ is the fluid-frame radiation flux. Thus the transport opacity
receives $\lambda_E/c$, in addition to ordinary absorption and scattering.
GRACE stores this pair/decay transport contribution separately to avoid
counting it again in the ordinary collision source.

\textbf{Equilibrium.} Thermal pairs vanish for FD spectra at the matter
temperature with $\mu_\nu+\mu_{\bar\nu}=0$. This is a family of equilibria.
Pairs alone do not select nucleon beta equilibrium.

\newpage
\section{Muon decay and inverse decay}
Both charge-conjugate reactions are included:
\begin{equation}
 \mu^-\rightleftharpoons e^-+\nu_\mu+\bar\nu_e,\qquad
 \mu^+\rightleftharpoons e^++\bar\nu_\mu+\nu_e.
\end{equation}
Use $\sigma=+1$ for $\mu^-$ and $-1$ for $\mu^+$. Daughter 0 always has muon
flavour, daughter 1 electron flavour. Define
\begin{equation}
 f_{\ell,\sigma}(E)=\left[e^{(E-\sigma\mu_\ell)/T}+1\right]^{-1},\qquad
 \Delta_\sigma=\sigma(\mueff-\mu_e).
\end{equation}
Normally $\mueff=\mu_\mu^{\rm EOS}$. Section 3 gives the optional modification.

\subsection{Optimized spectral approach}
The finite-mass kernel integrates the tree-level amplitude
\begin{equation}
 \overline{|\mathcal M|^2}=64G_F^2(p_e\!\cdot q_0)(p_\mu\!\cdot q_1)
\end{equation}
over charged-lepton phase space and neutrino relative angle. Its matter factors are
$f_\mu(E_e+\epsilon+\epsilon')[1-f_e(E_e)]$ for production and
$f_e(E_e)[1-f_\mu(E_e+\epsilon+\epsilon')]$ for absorption.
The positive sum $R^{\rm tot}=R^p+R^a$ is integrated first, then split stably:
\begin{equation}
 x=\frac{\epsilon+\epsilon'-\Delta_\sigma}{T},\qquad
 R^p=\frac{R^{\rm tot}}{1+e^x},\qquad
 R^a=\frac{R^{\rm tot}}{1+e^{-x}}.
 \label{eq:balance}
\end{equation}
Bounded exponentials avoid overflow and preserve the inverse channel when
forward decay is extremely small. Equations \eqref{eq:fit} and
\eqref{eq:event} apply to each cross-flavour daughter pair.
All four lepton-carrying neutrino species are solved together with ordinary
rates and muonic thermal pairs.

\textbf{What is optimized?} \code{leptonic_block_jacobian: true} uses analytic
occupation derivatives during Newton Jacobian probes. For example,
\begin{equation}
 \frac{\partial I_{ij}}{\partial f_i}
 =-w_iw_j\left[R^p_{ij}(1-f_j')+R^a_{ij}f_j'\right].
\end{equation}
The M1 closure and FD fit retain numerical probes. Acceptance uses the full
nonlinear residual, with the original numerical solver as fallback.
Optional \code{muon_decay_kernel_cache} and \code{thermal_pair_kernel_cache}
reuse kernels only for exactly matching matter and grid inputs on CPU.
These optimizations preserve the collision model and equilibrium conditions.

\subsection{Lifetime-based \texttt{grey\_decay} approach}
A 16-point Fermi-gas integral gives each charged-muon density $n_{\mu,\sigma}$
and mean total energy $U_\sigma$. With $M=m_\mu$ and
$\tau_\mu=2.197\,\mu\mathrm{s}$, the representative energies are
\begin{equation}
 \epsilon_0=0.35U_\sigma,\quad \epsilon_1=0.30U_\sigma,\quad
 E_e=0.35U_\sigma,\qquad
 A_0=\frac{n_{\mu,\sigma}}{f_{\mu,\sigma}(U_\sigma)}\frac{M}{U_\sigma\tau_\mu}.
\end{equation}
Writing $h=1-f$, the forward and inverse coefficients and event rate are
\begin{align}
 P&=A_0 f_\mu(U_\sigma)h_e(E_e),&
 A&=A_0 h_\mu(U_\sigma)f_e(E_e),\\
 r_\sigma&=P h_0(\epsilon_0)h_1(\epsilon_1)-A f_0(\epsilon_0)f_1(\epsilon_1),&
 Q_a&=\epsilon_a r_\sigma.
\end{align}
Coefficients are computed logarithmically. The daughters use the actual fitted
radiation spectra. $P/A=e^{-(\epsilon_0+\epsilon_1-\Delta_\sigma)/T}$ ensures
chemical detailed balance. No spectral decay matrix or spectral fallback is used.
The common-event solve and conservative feedback remain active. This is an
experimental representative-energy approximation. Archived wider-domain tests
found substantial rate errors and some wrong reaction signs relative to spectral
decay, despite good energy conservation. Detailed balance does not establish
the correct relaxation rate. The old spectral-interpolation prototype is obsolete.

\newpage
\section{Low \texorpdfstring{$Y_\mu$}{Ymu}, atmosphere and chemical potentials}
\subsection{Thermodynamic consistency of reversible decay}
Spectral and grey decay require the additive leptonic EOS with
\code{dilute_muon_suppression: false}. In this configuration,
\begin{equation}
 P=P_b(\rho,T,Y_e+Y_\mu)+P_e(\rho,T,Y_e)+P_\mu(\rho,T,Y_\mu),
\end{equation}
with analogous sums for energy and entropy. The charged-lepton tables supply
the chemical potentials. A small \emph{net} $Y_\mu$ does not imply that both
thermal muon populations vanish.

The historical suppression option instead sets muon thermodynamic contributions
to zero at $Y_\mu\leq6\times10^{-4}$ and ramps $\mu_\mu$ from zero over
$\Delta Y_\mu=5\times10^{-5}$. It belongs to the older vacuum-decay setup.
Combining that suppressed bath with reversible spectral/grey decay is rejected.
The old charged-muon rate gate is also disabled for these reversible modes.
Thermal neutrino-pair kernels always bypass that abundance gate.

\subsection{Optional dilute charged-current target correction}
The current opt-in fix is \code{m1.eas.dilute_muon_cc_target: true}.
Let $\mu_\mu^{\rm raw}$ denote the unsuppressed table value, $Y_0=6\times10^{-4}$,
and $dY=5\times10^{-5}$. The implemented reaction potential is exactly
\begin{equation}
 \boxed{\mueff=\begin{cases}
 \mu_\mu^{\rm raw},&\mu_\mu^{\rm raw}<\mu_e,\\[2pt]
 \mu_e+(\mu_\mu^{\rm raw}-\mu_e)
 \tanh\!\left[\dfrac{\max(Y_\mu-Y_0,0)}{dY}\right],
 &\mu_\mu^{\rm raw}\geq\mu_e.
 \end{cases}}
\end{equation}
On the demuonizing branch this approaches $\mu_e$ at the dilute floor, removing
the positive offset $\mu_\mu-\mu_e$ there. It modifies the WeakHub muonic
charged-current emissivity target while retaining the existing grey absorption
averages. The \emph{same} $\mueff$ enters spectral and grey decay and their
muon-neutrino equilibrium target. Otherwise charged-current reactions and decay
would impose inconsistent chemical affinities.

This is a phenomenological reaction-target regularization. It does not replace
the EOS pressure or energy, and it is not an exact spectral reintegration of
the WeakHub table. It requires a leptonic EOS, a noPP WeakHub table and
\code{betaeq_policy: off}. It defaults to false. There is no explicit density
factor in this formula, and it does not guarantee zero muon production at low
$\rho,Y_\mu$. Inverse decay remains possible when both incoming neutrino
partners are present. The branch $\mu_\mu^{\rm raw}<\mu_e$ is unchanged.

\subsection{Atmosphere and numerical safeguards}
For
\begin{equation}
 \rho<\rho_{\rm fl}(1+\mathrm{atmo\_tol}),
\end{equation}
EAS preparation skips the EOS and rate calculation, clears pair/decay activation
and transport additions, and writes the ordinary numerical rate floors.
This prevents active pair/decay microphysics in the designated atmosphere.
It does not impose an additional cutoff throughout all low-density matter.

Reversible modes restore the unclamped matter-based neutrino chemical potentials.
They do not use the historical muon-neutrino degeneracy clamp to $\pm5$ or
optical-depth suppression of those potentials. Stable logistic splitting,
direct evaluation of Fermi holes and scaled dilute LTE ratios address numerical
underflow without independently clipping the chemical affinities.
EOS composition floors remain. A shared acceptance factor limits radiation and
matter increments together when a composition or energy bound is reached.
Optional \code{fermi_fofc} limits inadmissible transported radiation states. It
is a realizability safeguard, not a switch suppressing muonization.

\newpage
\section{Beta equilibrium and conservative source update}
\subsection{Which equilibrium is imposed?}
All kernel pair modes require \code{m1.eas.betaeq_policy: off}.
The chemical, timescale and Gieg-style beta-equilibrium remappings are therefore
not applied. Physical equilibration proceeds through the ordinary charged-current
sources and the pair/decay sources with matter backreaction.

In GRACE's EOS convention, the equilibrium targets are
\begin{align}
 \mu_{\nu_e}^{\rm eq}&=\mu_e+\mu_p-\mu_n-Q_{np},\\
 \mu_{\nu_\mu}^{\rm eq}&=\mueff+\mu_p-\mu_n-Q_{np},\qquad
 \mu_{\bar\nu_\ell}^{\rm eq}=-\mu_{\nu_\ell}^{\rm eq}.
\end{align}
$Q_{np}$ is the neutron--proton mass difference in the code's table convention.
These are ordinary-rate equilibrium targets, not prescribed chemical potentials
for the evolving FD fits. The targets obey
\begin{equation}
 \mu_{\nu_\mu}^{\rm eq}+\mu_{\bar\nu_e}^{\rm eq}=\mueff-\mu_e,
 \qquad
 \mu_{\bar\nu_\mu}^{\rm eq}+\mu_{\nu_e}^{\rm eq}=-(\mueff-\mu_e),
\end{equation}
which are precisely the decay balance conditions in both spectral and grey
decay. With the dilute correction enabled, equilibrium refers to this modified
reaction potential. At transparent beta equilibrium the neutrino potentials
vanish. In trapped matter they can be nonzero and depend on conserved lepton
content. Decay alone enforces its cross-flavour chemical relation and does
not determine nucleon beta equilibrium.

\subsection{How the source step approaches equilibrium}
The four species each carry $(E,F_x,F_y,F_z,N)$, giving the connected equation
\begin{equation}
 U^{+}-U^*=h\left[S_{\rm ordinary}(U^{+})+
 S_{\rm thermal}(U^{+})+S_{\rm decay}(U^{+})\right].
\end{equation}
$h$ is the diagonal IMEX stage coefficient times the timestep. Matter and
kernel inputs are fixed within this solve. Thermal tau pairs are updated
separately. Without decay, each resolved thermal pair uses a ten-variable solve.
Full residual checks, positivity, flux-cone and FD reconstruction checks remain
active in the optimized implementation. Small number defects are projected onto
reaction stoichiometry after removing the ordinary contribution, followed by
another full residual check.

For decay alone, its common rates imply
\begin{equation}
 n_b\dot Y_e=r_- -r_+,\qquad
 n_b\dot Y_\mu=-r_-+r_+,\qquad \dot Y_e+\dot Y_\mu=0.
\end{equation}
In the code, feedback uses the actual radiation increments, including ordinary
reactions. With $\Delta$ denoting new minus old densitized quantities,
\begin{align}
 \Delta\tau_{\rm matter}&=-\sum_s\Delta E_s,&
 \Delta S_{i,\rm matter}&=-\sum_s\Delta F_{i,s},\\
 \Delta(DY_e)&=-\Delta N_{\nu_e}+\Delta N_{\bar\nu_e},&
 \Delta(DY_\mu)&=-\Delta N_{\nu_\mu}+\Delta N_{\bar\nu_\mu}.
\end{align}
$D$ is conserved baryon mass density and GRACE's $N$ includes the baryon-mass
normalization. One common $0\leq\theta\leq1$ scales the connected matter and
radiation increments at bounds. No separate $(m_\mu-m_e)r$ heating term is added,
because the leptonic EOS already contains the charged-lepton rest energy.
At $\theta<1$, exchange remains conservative but the accepted step is limited.
This fixed-matter stage is not a simultaneous matter--radiation equilibrium solve
at arbitrarily large stiffness.

\vfill
{\footnotesize\raggedright\textbf{Source map and scope.}
Checked against the local code at the commit on page 1, not an inferred run
configuration. Under \code{include/grace/physics/}: \code{neutrino_pair_collision.hh},
\code{neutrino_muon_decay.hh}, \code{neutrino_grey_decay.hh},
\code{neutrino_leptonic_update.hh}, \code{eas_policies.hh},
\code{eas_neutrino_rates_analytic.hh}, \code{eas_kinds.hh}, \code{m1.hh}
and \code{eos/leptonic_eos_4d.hh}. Defaults: \code{parameters/m1.yaml}.
Further derivations: \code{doc/developer_guide/pair_processes.md} and
\code{muon_decay.rst}. Grey-model comparison:
\code{calea_hottov_comparison/grey_domain_status.rst} in the parent workspace.
No simulation or new rate validation was performed for this note.\par}
\end{document}
